\section{Conclusion and Perspectives}
The three-tier framework makes placement, cloud scheduling and coordination
measurable under paired task replays. On the reference topology,
centralized DQN reaches 77.41\% deadline success, and the additional
centralized Double DQN control reaches 77.40\%. Greedy/EDF achieves
the lowest mean P95, at 12.95 s. Among hierarchical controllers,
Double DQN reduces the reference P95 from 16.91 to 15.08 s relative
to DQN, with a 0.30 percentage point decrease in deadline success.
The latency ranking reverses under combined stress.

The frozen-policy ablation identifies SPT as the strongest cloud rule
on the tested workloads. On the reference, replacing the learned
Double DQN cloud scheduler by SPT reduces mean P95 from 15.08 to
12.36 s with the same upstream weights. The event-triggered feedback
experiment also identifies a practical control mechanism. At
threshold 0.50, Double DQN transmits about 70\% fewer feedback bytes
than with frequent updates, with a 2.72 s increase in mean P95.

These results characterize one chronological test window and a
single-worker cloud. An independent trace period, multi-worker
execution and a shared data/control network would extend the
evaluation. Federated training is a further direction. Devices could
aggregate local updates at each edge before the edges synchronize at
the cloud, with synchronization traffic measured separately from
inference feedback.
